\documentclass{article}

\usepackage{amsmath,amssymb}
\usepackage{xurl}
\usepackage[hidelinks]{hyperref}

% Shared commands for notes in docs/paper-gaps/.

\DeclareUnicodeCharacter{2080}{\ifmmode{}_{0}\else\textsubscript{0}\fi}
\DeclareUnicodeCharacter{2081}{\ifmmode{}_{1}\else\textsubscript{1}\fi}
\DeclareUnicodeCharacter{2082}{\ifmmode{}_{2}\else\textsubscript{2}\fi}
\DeclareUnicodeCharacter{2099}{\ifmmode{}_{n}\else\textsubscript{n}\fi}

\newcommand{\Lean}{\textsc{Lean}}
\ExplSyntaxOn
\NewDocumentCommand{\leanid}{m}{
  \ifmmode
    \textsf{\detokenize{#1}}
  \else
    \group_begin:
    \urlstyle{sf}
    \tl_set:Nn \l_tmpa_tl {#1}
    \tl_replace_all:Nnn \l_tmpa_tl {\_} {_}
    \exp_args:NV \path \l_tmpa_tl
    \group_end:
  \fi
}
\ExplSyntaxOff
\providecommand{\tr}{\operatorname{tr}}
\newcommand{\tnleanIssueBase}{https://github.com/LionSR/TNLean/issues/}
\newcommand{\tnleanPrBase}{https://github.com/LionSR/TNLean/pull/}
\newcommand{\tnleanDocBase}{https://sirui-lu.com/TNLean/docs/}
\newcommand{\ghissue}[1]{\href{\tnleanIssueBase#1}{GitHub issue \##1}}
\newcommand{\ghpr}[1]{\href{\tnleanPrBase#1}{PR \##1}}
\newcommand{\doclink}[1]{\href{\tnleanDocBase#1.html}{\path{#1}}}

% Dirac notation and the identity operator, as in blueprint/src/macros/common.tex.
% \providecommand keeps note-local definitions authoritative.
\usepackage{dsfont}
\providecommand{\ket}[1]{|#1\rangle}
\providecommand{\bra}[1]{\langle #1|}
\providecommand{\braket}[2]{\langle #1 | #2 \rangle}
\providecommand{\ketbra}[2]{|#1\rangle\!\langle #2|}
\providecommand{\Id}{\mathds{1}}
\providecommand{\one}{\mathds{1}}
\providecommand{\C}{\mathbb{C}}
\providecommand{\MN}[1]{M_{#1}(\C)}
\providecommand{\MD}{M_D(\C)}

% External referencing for paper sources.  The build helper writes sanitized
% auxiliary files under build/paper-aux/ and excludes duplicated source labels.
\usepackage{xr}

% Import a generated paper auxiliary file with a caller-supplied label prefix.
% The candidate paths support builds from docs/paper-gaps/, the repository root,
% and build/paper-aux/ itself.
\newcommand{\importpaperaux}[2]{%
  \IfFileExists{../../build/paper-aux/#2.aux}{%
    \externaldocument[#1]{../../build/paper-aux/#2}%
  }{%
    \IfFileExists{build/paper-aux/#2.aux}{%
      \externaldocument[#1]{build/paper-aux/#2}%
    }{%
      \IfFileExists{#2.aux}{%
        \externaldocument[#1]{#2}%
      }{}%
    }%
  }%
}

% General external reference macros
\newcommand{\paperref}[2]{\ref{#1-#2}}
\newcommand{\papereqref}[2]{(\ref{#1-#2})}

% CPSV16 source (1606.00608 / MPDO-22-12-17-2.tex) external document and shortcuts
\importpaperaux{cpsv16-}{1606.00608/MPDO-22-12-17-2-tnlean}
\newcommand{\cpsvref}[1]{\paperref{cpsv16}{#1}}
\newcommand{\cpsveqref}[1]{\papereqref{cpsv16}{#1}}

% Verdict marker: \gapnote{<kind>}{<status>}. Kind is the policy
% classification (clarification, local-correction, scope-restriction,
% unfaithful, false-source, open-gap); status is open, wip (elimination
% underway), resolved, or historical. Machine-read by the site index;
% prints a verdict line.
\newcommand{\gapnote}[2]{\par\noindent\textbf{Verdict:} #1 (#2).\par}


\title{Wolf Proposition~1.5: Factor and Coordinate Direct-Sum Scope Restrictions}
\author{}
\date{2026-07-18}

\begin{document}
\maketitle

\section{The source statement}

Wolf's Proposition~1.5 classifies positive linear retractions onto an arbitrary
unital finite-dimensional $*$-subalgebra
\cite[Proposition~1.5 and Equation~(1.40)]{Wolf2012Quantum}.  After unitary
conjugation, the subalgebra has the form
\[
    \bigoplus_{k=1}^{K} M_{d_k}(\mathbb C)\otimes\mathbf 1_{m_k},
\]
and the proposition asserts that the retraction is a direct sum of weighted
partial traces, one for each summand.  The proof first removes inputs between
distinct summands and outputs in the wrong summand, and then proves the
weighted-partial-trace formula separately on each factor.  The corresponding
local source is
\path{Notes/WolfNoteTexSource/ch01_deconstructing_quantum.tex}, lines~536--570.

\section{Formalized component}

The theorem
\leanid{Matrix.exists_density_of_positive_retraction_onto_right_factor} proves
the complete factor argument for a positive retraction whose range is
\[
    \mathbf 1_m\otimes M_d(\mathbb C).
\]
It introduces no hypothesis absent from the one-factor specialization of the
source proposition.  In particular, the representing matrix is only required
to be positive semidefinite and to have trace one.  Thus the theorem is a
faithful scope-restricted case, but it is not the full direct-sum
classification of Proposition~1.5.

The theorem
\leanid{Matrix.coordinateDirectSumConditionalExpectation_isKrausCPTP_and_isConditionalExpectation}
constructs the normalized trace-preserving choice simultaneously on all
summands in displayed direct-sum coordinates.  In the tensor-factor order used
by TNLean, its range is
\[
    \bigoplus_{k=1}^{K}\mathbf 1_{m_k}\otimes M_{d_k}(\mathbb C),
\]
and its formula is obtained from Wolf's order by interchanging the two tensor
factors.  This theorem includes the coordinatewise partial traces and removal
of off-diagonal blocks, but it neither conjugates an arbitrary represented
subalgebra into these coordinates nor corestricts the ambient matrix-valued map
to that subalgebra.  It is therefore a coordinate construction, not the full
unitarily transported statement of Proposition~1.5.

\section{Resolution of the direct-sum argument}

The theorem
\leanid{StarSubalgebra.exists_block_densities_of_positive_retraction} combines
the one-factor theorem with the finite-dimensional block description.  Its
proof:
\begin{enumerate}
    \item conjugate and reindex the ambient matrices by the unitary block data;
    \item prove that the positive retraction annihilates inputs between
          distinct central summands;
    \item prove that an input supported in one summand has no output in another
          summand;
    \item apply the one-factor theorem to each diagonal summand; and
    \item conjugate the resulting direct sum of weighted partial traces back to
          the original subalgebra.
\end{enumerate}
The second and third items are the off-diagonal positivity arguments in Wolf's
proof preceding Equation~(1.42).  They are now proved using the complementary
central projection of each summand.  The remaining distinction between the
classification of a given retraction and the later construction of a
fixed-point projection is recorded in
\path{docs/paper-gaps/wolf_theorem6_14_fixed_point_projection_gap.tex}.

\bibliographystyle{alpha}
\bibliography{references}

\end{document}
